\chapter{配置点策略与预算分配}
\section{基线及公平比较}
固定隐藏特征后，配置点决定哪些残差和约束被经验目标看到。本文比较A固定完整矩形网格、B固定随机点、C内部PDE残差增量点及D区域与界面联合指标。A/B一次求解；C/D先以共同初始点求解，再在同一候选池中新增点并重求输出系数，隐藏参数保持不变。

在同一随机种子内，特征流、点组流、候选流和抽样流分别管理。B的最终随机点包含C/D初始点前缀；A使用规则网格，不宣称与随机初始集相同。所有方法使用同一外边界点、固定权重和最终测试点。五种子为11、23、37、51、71，所有运行及失败状态均保存，不能只取最好一次。

界面一个物理点提供值与通量两行，故
\begin{equation}
m_{\rm eq}=m_{I,1}+m_{I,2}+2m_\Gamma+m_B.
\end{equation}
固定物理点预算时D可改变方程数；固定方程预算时D可改变物理点数。本文分别试验并同时报告两种计数，避免把新增界面方程的收益误归因于选点位置。

\section{残差指标和候选抽样}
给定当前解，内部残差$r_i=-\beta_i\Delta u_i^N-f_i$，界面残差$r_u=[u^N]-g_u$、$r_q=[\beta\partial_nu^N]-g_q$。在固定尺度下定义
\begin{equation}\label{eq:indicators}
\eta_i=\operatorname{RMS}(r_i/\beta_i),\qquad
\eta_\Gamma=\operatorname{RMS}\sqrt{r_u^2+(r_q/\beta_{\max})^2}.
\end{equation}
这些是采样指标而非已证明的后验解误差估计器。高对比度归一化可能弱化某些物理量的重要性，故实验仍报告未归一化通量。

C按残差分布选择内部新增点，但组预算和界面新增数固定；D根据三个指标调整预算，再按各组位置残差选点。可用候选点上的概率为
\begin{equation}
p_\ell=\frac{.2}{M}+.8\frac{|r_\ell|}{\sum_j|r_j|},
\end{equation}
界面用式\eqref{eq:indicators}中相应点指标代替$|r_\ell|$。全零时退回均匀概率，抽样无放回。这种探索混合避免只集中在少数候选点，但不保证连续区域覆盖或最佳配置。

非均匀点改变经验目标所偏重的位置。本文没有逆密度权重，因此不声称近似同一个固定均匀积分；连续期望可写为$\int p(x)r(x)^2\,dx$，与均匀残差范数不同。残差概率及增量/替换采样思想已有先行研究\cite{rar}；本文只研究一种简单联合预算机制的实际表现。

\section{最低配额及整数预算}
每轮最低新增点$(16,16,8)$，成本向量在点预算下为$(1,1,1)$，在方程预算下为$(1,1,2)$。扣除最低配额后，剩余成本按$\eta_k+\varepsilon$比例分配，除以各组点成本得到实数配额，先取整，再以剩余成本和最大欠配额补齐。成本1组保证此实验中预算能用尽。

初始内部点128/侧、界面32。静态增量为$(80,80,32)$，对应每轮192点或224标量行。主平滑比较三轮，A/B/C最终内部368/侧、界面128；峰值比较增加六轮设置，并固定最终总行1232或1904，外边界240行不变。

完整流程为：\begin{enumerate}
\item 固定模型、特征、边界、预算、候选池及随机流，保存配置。
\item 使用初始点求解；在候选池计算当前残差，不使用真解误差选点。
\item 计算组指标及整数配额，按探索混合概率新增未使用点。
\item 重新组装并求输出系数，重复预设轮数，不按最终测试误差挑轮次。
\item 在从未参与选点的网格评估全部指标，保留配额轨迹、系数及完整时间。
\end{enumerate}

\section{消融、候选偏差及复杂度}
去PDE或去界面指标时，对应组指标不参与动态预算，位置抽样退回均匀，最低配额和总预算仍不变。这个消融会同时改变配额与位置，因此后续固定比例扫描是必要的混杂因素对照，不能把消融变好直接解释为某指标无价值。

原D下一轮组RMS仍包含已经选入训练的候选点，可能乐观。D\_unused仅排除这些点后计算组RMS；位置选择仍无放回。该未用池也受历史选择影响，不是独立同分布最终测试集。本文没有据此使用iid采样的理论保证。

总特征$K$、矩阵行$m$、候选点$M$时，固定维度下组装约$O(mK)$，稠密SVD在$m\ge K$时主要为$O(mK^2)$，候选评估约$O(MK)$。$T$轮增量方案包含$T+1$次求解及$T$次候选评估，必须与静态的一次求解比较完整成本。重复奇异值诊断也纳入本项目计时；不能只展示最终一次求解的时间。
